\section*{Sets and Indices}

There are no nontrivial sets or indices in the implemented model. The transportation modes are represented explicitly:
\[
\{\text{horse},\text{bicycle},\text{handcart}\}.
\]

\section*{Parameters}

All parameters are scalars read from \texttt{general\_parameters.csv}:

\begin{itemize}
    \item $P^{\text{req}}$ (\texttt{min\_produce}): minimum total produce that must be transported.
    \item $P^{\text{target}}$ (\texttt{produce}): stated produce-to-transport amount. This value is loaded by the code but not used in the optimization model.
    \item $u_h$ (\texttt{horse\_cap}): produce capacity per horse trip.
    \item $u_b$ (\texttt{bike\_cap}): produce capacity per bicycle trip.
    \item $u_c$ (\texttt{cart\_cap}): produce capacity per handcart trip.
    \item $e_h$ (\texttt{horse\_poll}): pollution per horse trip.
    \item $e_b$ (\texttt{bike\_poll}): pollution per bicycle trip.
    \item $e_c$ (\texttt{cart\_poll}): pollution per handcart trip.
    \item $E^{\max}$ (\texttt{max\_poll}): maximum allowable total pollution.
    \item $L_h$ (\texttt{min\_horse}): minimum required number of horse trips.
    \item $M$ (\texttt{M}): big-\(M\) constant used to enforce the either/or choice between bicycle and handcart; in the code, \(M=10000\).
\end{itemize}

For Instance 63, the loaded values are:
\[
P^{\text{target}}=1000,\quad P^{\text{req}}=1000,\quad
u_h=55,\quad u_b=30,\quad u_c=40,
\]
\[
e_h=80,\quad e_b=0,\quad e_c=0,\quad
E^{\max}=1000,\quad L_h=8,\quad M=10000.
\]

\section*{Decision Variables}

\begin{itemize}
    \item $x_h$ (code: \texttt{x\_horse}): number of horse trips, integer and nonnegative.
    \item $x_b$ (code: \texttt{x\_bike}): number of bicycle trips, integer and nonnegative.
    \item $x_c$ (code: \texttt{x\_cart}): number of handcart trips, integer and nonnegative.
    \item $y$ (code: \texttt{y} / \texttt{choose\_bike}): binary selector, where \(y=1\) allows bicycle trips and forbids handcart trips, while \(y=0\) allows handcart trips and forbids bicycle trips.
\end{itemize}

\section*{Objective}

Minimize total pollution:
\[
\min \; e_h x_h + e_b x_b + e_c x_c.
\]

\section*{Constraints}

Minimum transported produce:
\[
u_h x_h + u_b x_b + u_c x_c \ge P^{\text{req}}.
\]

Maximum total pollution:
\[
e_h x_h + e_b x_b + e_c x_c \le E^{\max}.
\]

Minimum number of horse trips:
\[
x_h \ge L_h.
\]

Exclusive choice between bicycle and handcart via big-\(M\):
\[
x_b \le My,
\]
\[
x_c \le M(1-y).
\]

Integrality and binary restrictions:
\[
x_h, x_b, x_c \in \mathbb{Z}_{\ge 0},
\qquad
y \in \{0,1\}.
\]

\section*{Notes on Implementation}

The model is implemented as a mixed-integer linear program and solved with \texttt{GUROBI\_CMD} through a PuLP-compatible interface.

The formulation should be interpreted exactly as coded:

\begin{itemize}
    \item The optimization uses \texttt{min\_produce} as the required transported amount, not \texttt{produce\_to\_transport}. Although both are loaded, only \texttt{min\_produce} appears in the constraints.
    \item The exclusivity requirement is modeled as ``at most one of bicycle and handcart can be used,'' enforced by the binary variable \(y\) and big-\(M\) bounds. The model does not require that either bicycle or handcart be used; one of them may be zero and the other positive, or both may be zero if the other constraints permit it.
    \item Because \(e_b=e_c=0\) in this instance, the objective minimizes only horse-trip pollution, while the lower bound \(x_h \ge L_h\) and produce requirement determine how much of the remaining transport is assigned to the chosen zero-pollution mode.
    \item No upper bounds are placed directly on trip counts other than those implied by the pollution limit, produce requirement, and the big-\(M\) logic.
\end{itemize}
